\documentclass[11pt]{article}
\usepackage[margin=1in]{geometry}
\usepackage{amsmath,amssymb,amsthm,mathtools,mathrsfs}
\usepackage{booktabs,array,enumitem}
\usepackage[T1]{fontenc}
\usepackage{lmodern,microtype}
\usepackage{xcolor}
\usepackage{hyperref}
\hypersetup{colorlinks=true,linkcolor=blue!40!black,citecolor=blue!40!black,
  urlcolor=blue!50!black,
  pdftitle={A quantitative refinement of the seven-eighths zero-free half-plane},
  pdfauthor={Zhongpai Gao}}

\newtheorem{theorem}{Theorem}[section]
\newtheorem{proposition}[theorem]{Proposition}
\newtheorem{lemma}[theorem]{Lemma}
\newtheorem{corollary}[theorem]{Corollary}
\theoremstyle{definition}
\newtheorem{definition}[theorem]{Definition}
\newtheorem{remark}[theorem]{Remark}
\newcommand{\Q}{\mathbb Q}
\newcommand{\R}{\mathbb R}
\newcommand{\C}{\mathbb C}
\newcommand{\OO}{\mathcal O_F}
\newcommand{\N}{\mathrm N}
\newcommand{\Rea}{\operatorname{Re}}
\newcommand{\Imma}{\operatorname{Im}}
\newcommand{\one}{\mathbf1}
\newcommand{\norm}[1]{\left\lVert#1\right\rVert}

\title{A quantitative refinement of the seven-eighths\\zero-free half-plane}
\author{Zhongpai Gao\\\small gaozhongpai@gmail.com}
\date{7 October 2026}
\begin{document}
\maketitle

\begin{abstract}
We refine OpenAI's zero-free half-plane for all Dirichlet $L$-functions
and all finite-order Hecke $L$-functions over $\Q(\sqrt{-3})$.
The new boundary is $20999/24000=7/8-1/24000$, with principal poles
allowed. We change the averaging scales in its cubic-theta probe,
prove the corresponding low-energy estimate, and use an exact
two-variable endpoint calculation to bound the remaining rows.
An explicit quadratic certificate gives the required margin. The moment
estimates remain unchanged; the proof includes the residue, tail and
normalization costs at the new scales.
\end{abstract}

\noindent\textbf{2020 Mathematics Subject Classification.}
11M26 (primary); 11M06, 11L05 (secondary).

\noindent\textbf{Keywords.}
Dirichlet $L$-function, Hecke $L$-function, zero-free half-plane,
cubic theta function, marked prime factors.

\section{Introduction}\label{gex:introduction}

Throughout, $F=\Q(\sqrt{-3})$. OpenAI \cite{OAI26} proves that every
Dirichlet $L$-function and every finite-order Hecke $L$-function over
$F$ is zero-free in $\Rea s>7/8$. Its proof compares a cubic-theta
character sum with a Mellin signal containing $1/L$. We retain this
construction and its moment estimates, and improve the comparison
after changing the averaging scales.

\begin{theorem}\label{gex:main}
Put $\theta=20999/24000$. Every finite-order Hecke $L$-function over
$F$ and every Dirichlet $L$-function, including imprimitive
presentations, is zero-free in $\Rea s>\theta$, with principal poles
at $s=1$ excepted. In particular this holds for the Riemann zeta
function. Every zero $\rho$ with $0<\Rea\rho<1$ of one of these
$L$-functions satisfies
\[
 \frac{3001}{24000}\le\Rea\rho\le\frac{20999}{24000}.
\]
\end{theorem}

The boundary is independent of the conductor and height. Constants
and thresholds in the character-sum estimates may depend on the
fixed target.

The extra prime length lowers the direct estimate by $1/24000$
while leaving the principal exponent $s-11/16$ unchanged. It also
changes the estimate for the remaining rows. Retaining the two
parameters in that estimate gives a quadratic that is positive
throughout the parameter rectangle. The
resulting margin pays the change of scales and the recovery errors.

Sections~\ref{gex:source-section}--\ref{gex:low-recovery}
establish the low estimate. Section~\ref{gex:high-section} proves
the endpoint bound, and Section~\ref{gex:recovery-section} completes
the recovery and continuation. The last section explains the limit
of this parameter refinement.

The unchanged inputs are listed below. The new arguments are the
modified-scale low estimate, the endpoint certificate and the
admission of the principal Euler correction below $7/8$.

\begin{center}
\small
\begin{tabular*}{\linewidth}{@{\extracolsep{\fill}}ll@{}}
\toprule
Input in \cite{OAI26} & Use here\\
\midrule
Probe and Poisson; \S\S\,4--7, 16
 & Character sum and principal signal\\
\addlinespace[2pt]
Reflection and Gram; Lemma 14.3, Prop.\ 15.2
 & Low estimate at new scales\\
\addlinespace[2pt]
Inverse and plain moments; \S\S\,17--19
 & Witness counts at $\kappa=3/4$\\
\addlinespace[2pt]
Contours and continuation; \S\S\,2, 10, 20
 & Recovery and zero exclusion\\
\bottomrule
\end{tabular*}
\end{center}
